\usepackage[T2A]{fontenc}
\usepackage[utf8x]{inputenc}
\usepackage[russian]{babel}
\usepackage{fancyhdr}
\usepackage{amsmath}
\usepackage{amsthm}
\usepackage{amssymb}
\usepackage{changepage}
\usepackage{titlesec}
\usepackage{enumitem}
\usepackage{textcomp}
\usepackage{hyperref}
\usepackage{pgfplots}
\usepackage{indentfirst}
\usepackage{wrapfig}
\usepackage{xcolor}
\usepackage{etoolbox}

% --- настройка ссылок для оглавления ---
\hypersetup{
    colorlinks=true, %set true if you want colored links
    linktoc=all,     %set to all if you want both sections and subsections linked
    linkcolor=blue,  %choose some color if you want links to stand out
}
\setcounter{tocdepth}{1}
% --- настройка ссылок для оглавления ---

% --- настройка глав ---
\makeatletter
\patchcmd{\chapter}{\clearpage}{}{}{}
\makeatother

% --- KOMA-Script title fonts ---
\setkomafont{title}{\centering\bfseries\Huge}
\setkomafont{subtitle}{\centering\Large}
\setkomafont{author}{\centering\normalfont}
\setkomafont{date}{\centering\normalfont}

\renewcommand{\thechapter}{\Roman{chapter}}

% Заголовок главы: номер сверху, название ниже, всё по центру
\titleformat{\chapter}[display]
  {\normalfont\LARGE\bfseries\filcenter} % общий формат + выравнивание по центру
  {\chaptername~\thechapter}             % метка: "Глава I"
  {1ex}                                  % вертикальный зазор между меткой и названием
  {\Large}                               % формат самой строки названия
% при желании можно добавить [after-code] пятым аргументом

% Отступы вокруг заголовка
\titlespacing*{\chapter}{0pt}{3ex plus 1ex}{2ex}

\renewcommand{\thesection}{\arabic{section}}
\renewcommand{\thesubsection}{\thesection.\arabic{subsection}}      % или \thesection.\arabic{subsection}, если нужно 1.1
\renewcommand{\thesubsubsection}{\thesection.\thesubsection.\arabic{subsubsection}}

\titleformat{\section}
  {\normalfont\Large\bfseries} % Форматирование
  {§\thesection} % Нумерация
  {1em} % Расстояние между номером и заголовком
  {}

% Настройка для подсекций
\titleformat{\subsection}
  {\normalfont\bfseries} % Форматирование
  {\thesubsection} % Нумерация
  {1em} % Расстояние между номером и заголовком
  {}


% --- настройка глав ---

% --- Настройка колонтитула ---
\pagestyle{fancy}
\fancyhead{}\fancyfoot{}
\renewcommand{\sectionmark}[1]{\markboth{\thesection\ #1}{}}
\renewcommand{\subsectionmark}[1]{\markright{\thesubsection\ #1}}

\fancyhead[EOC]{\nouppercase{\leftmark}}
\fancyfoot[C]{\thepage}
\renewcommand{\headrulewidth}{0.4pt}
\renewcommand{\footrulewidth}{0.4pt}

\fancypagestyle{plain}{%
  \fancyhead{}\fancyfoot{}
  \fancyfoot[C]{\thepage}
  \renewcommand{\footrulewidth}{0.4pt}
  \renewcommand{\headrulewidth}{0pt}
  
}
% --- Настройка колонтитула ---

% --- Теоремы, определения ---
\theoremstyle{definition}
\newtheorem{theorem}{Теорема}
\newtheorem{definition}{Определение}
\newtheorem{lemma}{Лемма}
\newtheorem*{remark}{Замечание}

\newcounter{case}


% --- Теоремы, определения ---
